\documentclass[12pt]{article}
\usepackage[utf8]{inputenc}
\usepackage[T1]{fontenc}
\usepackage{hyperref}
\title{Yis Language Resources for K-12 Curriculum Development}
\author{Compiled by Claude Sonnet 5.5 (claude-sonnet-5-5) for Andr\'as Kornai}
\date{}
\begin{document}
\maketitle

\section{Introduction}
Yis (ISO 639-3 \texttt{yis}, Glottocode yiss1240) is a Torricelli language of Papua New Guinea. Wikipedia places it in Sandaun Province and Joshua Project places it in East Sepik Province, so sources disagree. Counts also differ: about 800 speakers (Wikipedia, 2014) against 600 (Joshua Project). According to Joshua Project, Yis is used for most daily talk and Tok Pisin with outsiders. Do not confuse it with Yeri, a separate Torricelli language.

\section{Primary Educational Resources}
None found.

\section{Linguistic and Reference Materials}
\begin{itemize}
\item Yis language, Wikipedia: \url{https://en.wikipedia.org/wiki/Yis_language}
\item Endangered Languages Catalogue entry: \url{https://endangeredlanguages.com/lang/10542?hl=ru}
\item Jennifer Wilson 2017, Glottolog reference entry (the search result did not show the title; consult it for the underlying work): \url{https://glottolog.org/resource/reference/id/552938}
\end{itemize}

\section{Religious and Community Resources}
\begin{itemize}
\item Joshua Project, Yis in Papua New Guinea (people-group profile; it reports no Bible portions translated): \url{https://joshuaproject.net/people_groups/16043}
\item Joshua Project language page: \url{https://joshuaproject.net/languages/yis}
\end{itemize}

\section{Digital and Audio Resources}
None found.

\section{Curriculum Development Context}
Yis has a small community that still uses the language daily. Documentation appears limited. The Wilson (2017) work should be located as a possible basis for an orthography and teaching materials. No school use was found.
\end{document}